% 11_variations.tex - ASCII grids, a missing :InputCRS, a refined grid
\section{Try three changes}
Make each change, run, and then undo the change before the next one. Each run writes to its own output folder.

\step{Use ESRI ASCII grids instead of GeoTIFF}
Change the two raster names in \cmd{Creek.rvg}:
\diff{t2_ascii.diff}
\begin{shell}
$ $RAVEN Creek -o out_ascii/
\end{shell}
\see \msg{Exiting Gracefully: Successful Simulation}.

\why the two pairs of rasters hold the same values (the GeoTIFF files store them as 32-bit numbers). The heads in
\cmd{out\_ascii/GWHRUState.csv} differ from those in \cmd{out/} by at most \V{t2:asciidiff}~m. An ESRI ASCII grid has no
projection of its own: \cmd{:InputCRS} applies to it, and Raven only checks that the grid covers the HRUs. Undo the change
(\cmd{dem.tif}, \cmd{bedrock.tif}) before the next step.

\step{Leave out \texttt{:InputCRS}}\label{st:nocrs}
Make the \cmd{:InputCRS} line a comment:
\diff{t2_nocrs.diff}
\term{t2_nocrs_err.txt}
\see Raven stops before it builds anything. Without \cmd{:InputCRS}, the coordinates of the well (and of every input file)
would be read as longitude and latitude. The message says what to do. (If the wells had been given by latitude and
longitude, the next message would have named \cmd{hrus.geojson}, whose \cmd{"crs"} member says EPSG:32611.) Undo the change
before the next step.

\step{Use smaller cells along the creek}
Add one line after \cmd{:GridCellSize}:
\diff{t2_refine.diff}
\begin{shell}
$ $RAVEN Creek -o out_refined/
\end{shell}
\term{t2_refine_summary.txt}
\why a regular grid is refined in whole rows and columns. The grid now has \V{t2:refnrow}$\times$\V{t2:refncol} cells and
\V{t2:refactive} active columns (instead of \V{t2:nrow}$\times$\V{t2:ncol} and \V{t2:activecols}), \V{t2:refriv} river
cells (instead of \V{t2:rivcells}), and a run takes about \V{t2:refsec}~seconds (instead of \V{t2:runsec}).
The smaller cells along the creek change the discharge to it: \V{t2:reftoriv} instead of
\V{t2:toriv}~m\textsuperscript{3}/d. The narrower river cells take less of the groundwater, and more of it seeps out into
the soil next to them (seepage to Raven \V{t2:refseep} instead of \V{t2:seep}~m\textsuperscript{3}/d). The mean
water-table depth of HRU~6 becomes \V{t2:refwt6}~m. Another option, \cmd{:GridType QUADTREE}, refines the grid locally
instead of in whole rows and columns. Undo the change at the end.
